\documentclass[12pt]{article}
\usepackage[utf8]{inputenc}
\usepackage[T1]{fontenc}
\usepackage{hyperref}
\title{Nong Zhuang Language Resources for K-12 Curriculum Development}
\author{Compiled by Claude Sonnet 5.5 (claude-sonnet-5-5) for Andr\'as Kornai}
\date{}
\begin{document}
\maketitle

\section{Introduction}
Nong Zhuang (ISO 639-3 zhn, Glottocode nong1247) is a Tai-Kadai (Kra-Dai) language of southern China, part of the Zhuang macrolanguage, with about 500,000 speakers. It is not known to be taught in schools. Web search found little material specific to Nong Zhuang. This report lists only what the searches returned.

\section{Primary Educational Resources}
No primers, textbooks or school materials specific to Nong Zhuang were found in search results.

\section{Linguistic and Reference Materials}
\begin{itemize}
\item Trilingual Nong Zhuang lexicon (Webonary), with more than four thousand words collected over about a decade, mostly from Zhetu District, Guangnan County. Its romanization follows the Chinese government-approved Zhuang orthography and its IPA transcriptions are based on the Nong dialect. Introduction: \url{https://www.webonary.work/zhuangwen/overview/introduction/}
\item Foreword to the same lexicon: \url{https://www.webonary.work/zhuangwen/overview/foreword/}
\item Wikipedia article on the Nong Zhuang language: \url{https://en.wikipedia.org/wiki/Nong_Zhuang_language}
\item Wikitongues language page: \url{https://wikitongues.org/languages/zhn/}
\end{itemize}
No standalone grammar was found in the search results.

\section{Religious and Community Resources}
\begin{itemize}
\item Joshua Project language page, which lists Scripture portions and a New Testament: \url{https://www.joshuaproject.net/index.php/languages/zhn}
\item Global Recordings Network page. It reported no recordings available in this language at the time of the search. \url{https://globalrecordings.net/en/language/zhn}
\end{itemize}

\section{Digital and Audio Resources}
The Webonary lexicon above is the main digital resource found. No audio corpus or software was found.

\section{Curriculum Development Context}
Nong Zhuang has no known school use. A curriculum would start from the Webonary lexicon and the government-approved Zhuang romanization it uses, and would need new primers and readers. Community consultation is needed first.

\end{document}